\chapter*{Conclusioni}
\chaptermark{Conclusioni}

L'obiettivo principale di questa tesi è stato lo studio degli eventi candidati muoni all'interno di JUNO per i primi dati raccolti dal rivelatore. I muoni, interagendo con il rivelatore, producono isotopi radioattivi che compongono il background cosmogenico; è quindi fondamentale lo sviluppo di algoritmi in grado di classificare e caratterizzare in maniera efficiente e precisa il tracciato dei muoni, in modo tale da applicare tecniche di veto. Ho esaminato dati provenienti da diverse \textit{Run}, utilizzando due metodi per la classificazione degli eventi: \textbf{correlazione temporale} e \textbf{clustering di carica} per i PMT della WP. I due algoritmi utilizzano approcci diversi per l'identificazione dei muoni e risultano quindi indipendenti.

Il primo metodo di selezione si basa sulla \textbf{correlazione temporale} tra eventi registrati nella WP e nel CD. A causa del passaggio rettilineo del muone nel rivelatore, si attende che il tempo di apertura delle finestre di trigger per l'acquisizione nella WP e CD siano molto ravvicinati temporalmente. Quindi, impostando un tempo di correlazione $t_\text{co}$ fisso, è stato possibile selezionare eventi candidati muoni che attraversano l'intero rivelatore (\ref{eq:diff}). Oltre al tempo, ho studiato lo spettro della carica degli eventi candidati, impostando due soglie ($E^\text{WP}_\text{th}$ e $E^\text{CD}_\text{th}$). In questa maniera ho ottenuto due rate degli eventi candidati muoni calcolati con il metodo di correlazione temporale per la RUN4049: $(4.04\pm0.02)\,\unit{Hz}$ e $(4.62\pm0.02)\,\unit{Hz}$ rispettivamente effettuando la selezione nello spettro di carica della WP e del CD.

Il secondo metodo si basa sul \textbf{clustering di carica} per i PMT della WP, ed è implementato nel software ufficiale di JUNO. Si basa sull'identificazione di "gruppi" di PMT vicini tra loro (cluster) attivati dal passaggio del muone. Il clustering consente di identificare anche muoni clipping che non attraversano direattamente il CD, poiché i PMT della WP sono posizionati su una struttura in acciaio inossidabile con un raggio maggiore rispetto al Central Detector. Utilizzando questo metodo per i dati provenienti dalla RUN4049 è stato ottenuto un rate  di eventi candidati muoni pari a $(5.66\pm0.02)\,\unit{Hz}$. Grazie ad altre informazioni legate alla posizione e alla direzione degli eventi candidati, è stato anche possibile studiare la distribuzione dell'angolo polare e azimutale ottenendo risultati attesi: gli eventi candidati muoni hanno una direzione preferibilmente verticale discendente e una distribuzione azimutale più complessa che dipende dalla morfologia del territorio circostante in accordo con simulazioni Monte Carlo.

Dall'analisi dei due metodi, è emerso che la frequenza degli eventi calcolati con il metodo di correlazione temporale risulta inferiore del $28.6\%$, considerando la selezione sullo spettro di carica della WP, e del $18.4\%$, considerando la selezione sullo spettro di carica del CD, rispetto al metodo di clustering di carica per i PMT della WP. Questa discrepanza può essere spiegata proprio dal diverso metodo di selezione degli eventi: il primo metodo seleziona eventi che sono stati rivelati anche nel CD, mentre il secondo metodo include anche i muoni clipping. Le simulazioni Monte Carlo e le estrapolazioni basate su misure ottenute in altri laboratori sotterranei suggeriscono che il valore atteso di muoni nel CD di JUNO si aggiri intorno a $\sim4\,\unit{Hz}$, confrontabile con i risultati ottenuti.

Da un lato, il metodo di correlazione temporale fornisce una selezione più affidabile, in quanto il criterio di correlazione temporale è più stringente e la probabilità di falsi positivi è inferiore rispetto al secondo metodo. D'altra parte, il primo metodo fornisce una quantità di informazioni più limitata, poiché considera solo il tempo assoluto e la carica di ciascun evento. Il secondo metodo, invece, offre una maggiore quantità di informazioni grazie alla ricostruzione del centro del cluster dei PMT e alla capacità di identificare anche i "muoni bundle".

In conclusione, sebbene lo studio dei muoni sia ancora in fase di ottimizzazione e i due algoritmi presentino solo una parte limitata della strategia di classificazione e caratterizzazione degli eventi candidati muoni, i risultati ottenuti evidenziano la difficoltà nel trovare un metodo definitivo di selezione. Un approccio combinato potrebbe quindi rappresentare una soluzione promettente per migliorare la selezione e caratterizzazione degli eventi candidati muoni nel rivelatore JUNO.
